\section{Review of Related Literature and Theoretical Framework}

\subsection{Safe Haven Currencies}

The currency literature builds on the hedge and safe haven definitions set out in the introduction. The safe haven concept in its modern form begins with the distinction between hedging and safety. \textcite{baur_lucey_2010} show for gold that an asset can be uncorrelated with equities and bonds on average without providing protection in a crash, and their hedge-safe haven dichotomy has become the standard vocabulary. \textcite{hossfeld_macdonald_2014} apply the distinction to exchange rates and find that the Swiss franc and the US dollar qualify as safe haven currencies among the G10, while yen appreciation in crises is driven mainly by the unwinding of carry trades rather than by safe haven inflows. The empirical record on the yen is otherwise strong. \textcite{ranaldo_soderlind_2010} find that the Swiss franc and the yen appreciate against the dollar and the euro when global equity markets fall and volatility rises. \textcite{fatum_yamamoto_2015} go further, ranking the yen as the safest of the safe haven currencies because it appreciates as market uncertainty increases regardless of the prevailing level of uncertainty, whereas other currencies react nonlinearly. \textcite{cho_han_2021} confirm in a cross-quantilogram framework that the yen, the euro, and the Swiss franc appreciate in periods of high foreign exchange volatility and rising US Treasury yields.

The literature also identifies the structural conditions that make a currency a safe haven. \textcite{habib_stracca_2012} stress net foreign asset positions, which survive as the most consistent predictor of safe haven status after controlling for carry trade, along with financial development and foreign exchange market liquidity. Low interest rates create the carry trade that positions the yen as a funding currency, and \textcite{nishigaki_2007} shows that US stock prices dominate the behavior of yen carry positions while the interest rate differential with Japan plays a smaller role. \textcite{brunnermeier_2009} demonstrate that carry trade returns crash when funding currencies appreciate, which is precisely the mechanism that produces yen strength during risk-off episodes. \textcite{imakubo_kamada_kan_2015} formalize the demand side through a currency premium that depends on deviations of real interest rates and real exchange rates from equilibrium.

\subsection{The Yen as a Safe Haven Currency}

Japan-specific evidence refines the general picture in one important respect. \textcite{botman_carvalho_lam_2013} conduct a forensic analysis of yen appreciation during risk-off episodes and find that the appreciation is unrelated to capital inflows and to expectations about relative monetary policy. The appreciation coincides instead with portfolio rebalancing through offshore derivative transactions, where reduced net short positions and a buildup of net long positions in the yen move the spot rate without appearing in the balance of payments. This finding matters for the present study because it predicts a null response of measured capital flows alongside a strong exchange rate response, exactly the pattern that \textcite{beirne_sugandi_2023} report.

The real economy consequences of safe haven status are well documented for Japan. \textcite{masujima_2017} argues that yen strength driven by safe haven demand slows the post-crisis recovery through net exports, and that large-scale monetary easing in response may mask financial vulnerabilities related to public debt sustainability. \textcite{belke_volz_2020} trace the appreciation to the structural hollowing out of Japanese industry. \textcite{iwaisako_nakata_2017} caution against overstating the export channel, noting that aggregate global demand played the larger role in the trade decline after the 2008 crisis, but their structural VAR still finds economically meaningful exchange rate effects on exports. On the bond side, \textcite{lam_tokuoka_2013} attribute the low and stable JGB yields to steady domestic inflows, high domestic ownership, and safe haven demand, while warning that population aging and the decline of private-sector savings threaten this stability.

\subsection{Real Effects and Financial Vulnerabilities}

Beyond the Japan-specific evidence reviewed above, the general literature on safe assets connects safe haven demand to the real economy through equilibrium interest rates and adjustment costs. \textcite{habib_stracca_venditti_2020} show that high demand for safe assets in crisis times depresses the natural real interest rate, with the exchange rate appreciation and the contraction in global demand disrupting normal adjustment mechanisms. \textcite{bussiere_lopez_tille_2013} document that large real appreciations create adjustment costs that can lead to economic dislocation when exchange rates eventually revert. For economies with low inflation and policy rates near the zero bound, real appreciations driven by risk-off episodes can feed deflationary pressures and weigh on aggregate demand. These mechanisms give the negative output response to risk-off shocks in safe haven destinations a clear theoretical basis.

Uncertainty expectations interact with the exchange rate through a separate channel. \textcite{beckmann_czudaj_2016} find that exchange rate expectation errors decrease when uncertainty rises, particularly for monetary policy uncertainty, and that the yen is exceptional in appreciating when uncertainty increases. This result is relevant to the uncertainty variable in the empirical model, because it suggests that the yen's response to risk-off shocks operates through a repricing of risk rather than through a measurable rise in domestic economic policy uncertainty, which would reconcile a muted uncertainty response with a strong currency response.

\subsection{Risk-Off Episodes and Capital Flows}

The mechanisms described above are triggered by identifiable episodes of market stress, and the empirical definition of a risk-off episode used throughout this literature follows \textcite{debock_carvalhofilho_2013}, who identify a risk-off event as a day on which the Chicago Board Options Exchange Volatility Index (VIX) exceeds its 60-day backward-looking moving average by 10 percentage points. Their study of currency behavior during such episodes finds that high-yield currencies depreciate and low-yield currencies appreciate, consistent with carry trade unwinding, and that the yen behaves as the clearest case of a low-yield currency gaining in stress periods.

Capital flow dynamics in response to risk-off shocks are governed by portfolio rebalancing. \textcite{hau_rey_2006} document a strong negative comovement between exchange rates and relative equity market returns, driven by repatriation and portfolio rebalancing flows. When global investors sell foreign equities and repatriate funds, the funding currency appreciates even without a change in measured portfolio inflows. This repatriation mechanism, together with the offshore derivative channel of \textcite{botman_carvalho_lam_2013}, explains why capital flows measured in the balance of payments can be unresponsive to risk-off shocks while the exchange rate moves sharply.

\subsection{The Post-2021 Regime and the Yen's Safe Haven Status}

The safe haven evidence summarized above is built almost entirely on data through 2021, and the years since then have called the yen's status into question. \textcite{beirne_sugandi_2023} themselves flag the onset of the break. In their footnote on the period following the Russian invasion of Ukraine in February 2022, they note that widening yield differentials between the United States and Japan, rising energy prices, and the worsening of Japan's balance of payments position drove a sharp yen depreciation against the dollar, with the dollar soaring against traditional safe haven currencies including the yen and the Swiss franc. The Federal Reserve's aggressive tightening cycle, which took the policy rate to 5.25 to 5.5 percent against a near-zero Bank of Japan policy rate, pushed the dollar-yen rate to 145 by September 2022 and past 150 by October, its weakest level since 1990. The Ministry of Finance intervened in September and October 2022, its first intervention since 1998, spending roughly 9.2 trillion yen to arrest the slide, and the currency then traded between 130 and 152 through 2023 as the Bank of Japan twice adjusted its yield curve control framework. The push to 160 in April 2024 drew a second intervention round of roughly 9.8 trillion yen, the first since October 2022, which capped the depreciation near that level. The yen's crisis appreciation became conditional on the rate differential rather than automatic.

Recent empirical work formalizes this conditionality. \textcite{park_fang_2025} estimate time-varying intra-safe haven currency behavior among the dollar, the franc, and the yen and find that the hierarchy among safe haven currencies shifts over time, with the yen's crisis response depending on the type of shock and the prevailing monetary policy constellation. \textcite{lu_yu_li_li_2024} reach the same conclusion from the East Asian market, showing that no East Asian currency, including the yen, qualifies as a safe haven when uncertainty is measured by geopolitical risk or trade policy uncertainty, while the yen retains safe haven properties under the conventional VIX-based risk measure. The International Monetary Fund's surveillance attributes the yen's persistent weakness to the interest rate differential with the United States and documents a 2024 depreciation of roughly 8 percent with exchange rate volatility beyond what rate expectations alone explain \parencite{imf_2024,imf_2025}. This stands in contrast to the unconditional ranking of \textcite{fatum_yamamoto_2015}, who found the yen to be the safest of the safe haven currencies over the global financial crisis period. The difference between the two studies is not a contradiction. It is evidence that the safe haven mechanism itself is regime-dependent, which is precisely the property that a sample extension can test.

The August 2024 episode sharpens the point. A Bank of Japan rate hike combined with a weak United States jobs report triggered a rapid unwind of leveraged yen carry positions in early August 2024, with the TOPIX index falling 12.4 percent on 5 August and the VIX surging above 65, a risk-off rally in the yen driven substantially by carry trade deleveraging rather than a vote of confidence in the Japanese economy \parencite{aquilina_lombardi_schrimpf_sushko_2024,bis_2024,imf_2025}. The yen appreciated sharply during this episode, but the appreciation was the mechanics of carry trade deleveraging rather than a fresh wave of safe haven inflows. The episode sits inside the extension window of the present study and appears in the risk-off episode table in Section 4.

Domestic policy developments complete the picture. The Bank of Japan ended yield curve control and exited its negative interest rate policy in March 2024, its first policy rate increase since 2007, and raised the policy rate in five steps to 1.00 percent by June 2026, its highest level in three decades, with 10-year JGB yields rising above 2.5 percent after more than a decade near zero, and the Nikkei 225 more than doubled between 2023 and 2026 on the back of corporate governance reforms and the return of inflation. The April 2025 tariff shock provided the clearest counterexample to the depreciation narrative, with the yen strengthening toward 140 as the dollar failed to behave as a safe haven, through carry trade unwinding and haven flows, before the currency weakened again toward 160 by mid-2026 as the rate differential reasserted itself. The post-2021 sample therefore contains a different monetary policy constellation, a different external position, and a different equity market regime. Whether the structural relationships estimated over the paper period survive in this environment is an open empirical question, and it is the question that the extended sample in this study is designed to answer.

\subsection{Theoretical Framework}

The theoretical framework synthesizes the evidence reviewed above into three transmission channels from a global risk-off shock to the Japanese economy. The first is flight to safety, through which a rise in global risk aversion increases demand for yen-denominated assets, appreciating the yen in real effective terms and compressing JGB yields relative to US Treasury yields. The second is carry trade unwinding, through which the yen, as the traditional funding currency, is bought back by leveraged investors when a risk-off shock hits, amplifying the appreciation beyond what safe haven demand alone would produce. The third is portfolio rebalancing and offshore derivative activity, through which investors adjust yen positions in derivatives and repatriation flows that move the exchange rate without registering as portfolio inflows in the balance of payments, separating the price response from the volume response.

The framework yields directional predictions for the endogenous variables in the model. A risk-off shock should appreciate the yen REER, widen the JGB yield spread relative to the United States, and depress real GDP through the competitiveness channel with a lag of one to two quarters. Capital flow variables should show no systematic response, because the price channels adjust through offshore activity. Domestic economic policy uncertainty should be muted relative to other safe havens, because market participants price in the repatriation of overseas holdings. The framework does not deliver an unambiguous prediction for the stock market, where the surge-to-safety effect on Japanese equities competes with the negative wealth and exporter effects of appreciation.

\subsection{Statement of Hypotheses}

To evaluate the five transmission channels established in the theoretical framework, five formal hypothesis pairs are formulated. Each hypothesis pair consists of a Null Hypothesis ($H_{0}$) of zero effect and an Alternative Hypothesis ($H_{1}$) derived from international macroeconomic theory

\begin{itemize}
    \item \textit{Hypothesis 1 (Exchange Rate Appreciation Channel)}
\newline $H_{0,1}$: A global risk-off shock has no statistically significant effect on Japan's real effective exchange rate.
\newline $H_{1,1}$: A global risk-off shock induces a statistically significant and persistent real effective appreciation of the Japanese yen.

    \item \textit{Hypothesis 2 (Real Output Contraction Channel)}
\newline $H_{0,2}$: A global risk-off shock has no statistically significant effect on Japan's real gross domestic product.
\newline $H_{1,2}$: A global risk-off shock induces a statistically significant decline in Japan's real gross domestic product through the real exchange rate competitiveness channel with a lag.

    \item \textit{Hypothesis 3 (Sovereign Yield Spread Widening Channel)}
\newline $H_{0,3}$: A global risk-off shock has no statistically significant effect on the yield spread between 10-year Japanese Government Bonds and 10-year US Treasury notes.
\newline $H_{1,3}$: A global risk-off shock induces a statistically significant widening of the yield spread as safe-haven demand compresses Japanese yields relative to US yields.

    \item \textit{Hypothesis 4 (Capital Flow Neutrality Channel)}
\newline $H_{0,4}$: A global risk-off shock has no statistically significant effect on net portfolio investment inflows to Japan.
\newline $H_{1,4}$: A global risk-off shock induces a statistically significant change in net portfolio investment inflows to Japan.
\newline \textit{Theoretical Expectation}: The null hypothesis $H_{0,4}$ will fail to be rejected because safe-haven exchange rate adjustment operates through offshore derivative rebalancing and unmeasured repatriation rather than balance of payments volume inflows.

    \item \textit{Hypothesis 5 (Domestic Policy Uncertainty Channel)}
\newline $H_{0,5}$: A global risk-off shock has no statistically significant effect on Japan's domestic economic policy uncertainty index.
\newline $H_{1,5}$: A global risk-off shock induces a statistically significant increase in Japan's domestic economic policy uncertainty index.
\newline \textit{Theoretical Expectation}: The null hypothesis $H_{0,5}$ will fail to be rejected because market participants price in the safe-haven repatriation mechanism without raising domestic policy uncertainty.
\end{itemize}

Section 3 details the two-tier structural vector autoregression data pipeline and econometric methodology designed to test these five hypotheses.
